\section{Additional Multi-Severity Analysis}
\label{sec:appendix-multiseverity}

\providecommand{\vTwoNumberMacros}{tables/multiseverity_numbers}
\input{\vTwoNumberMacros}
\providecommand{\vTwoAssociationTable}{tables/table_multiseverity_associations}
\providecommand{\vTwoDecisionTable}{tables/table_multiseverity_decisions}

This section gives the aggregation and sensitivity details behind
\Cref{sec:exp-multiseverity-extension}. The broad-severity experiment is
analyzed separately from the original narrow sweep. It keeps the original
thresholds and pair construction fixed and uses continuous retention rather
than defining a new performance band.

The same \vTwoVetoAtZeroEight{} TwoRoom checkpoints account for every
additional SR rejection among IR-passing rows at
$\sigma_{\mathrm{eval}}=\vTwoEvalZeroEight$ and
$\vTwoEvalZeroSixteen$: \vTwoLowVetoAtOne{} has
$\sigma_{\max}=\vTwoTrainOne$ and \vTwoLowVetoAtOneFive{} have
$\sigma_{\max}=\vTwoTrainOneFive$. At
$\sigma_{\mathrm{eval}}=\vTwoEvalZeroThirtyTwo$, these rejections comprise
\vTwoVetoThirtyTwoTwoRoom{} TwoRoom, \vTwoVetoThirtyTwoPushT{} PushT, and
\vTwoVetoThirtyTwoReacher{} Reacher checkpoints; at
$\vTwoEvalZeroSixtyFour$, they comprise
\vTwoVetoSixtyFourTwoRoom{} TwoRoom and \vTwoVetoSixtyFourPushT{} PushT
checkpoints. No Cube checkpoint passes IR while failing SR. The conditional
rejections thus
span three tasks at $\vTwoEvalZeroThirtyTwo$ and two at
$\vTwoEvalZeroSixtyFour$, although some checkpoint rows recur across severity
slices and the conditional denominator also changes. All low-severity
rejections fall on TwoRoom, whose SR catalog uses one median split and has the
fewest eligible anchors (61 of 100; \vTwoSRLabelReference); this concentration should be
read within the scope of the tested pair catalog.

\input{\vTwoDecisionTable}
\FloatBarrier

The reverse disagreement in
\Cref{tab:multiseverity-decision-decomposition} is expected to be possible
because the conditions use different references. Relative IR divides the raw
perturbation radius by the unaugmented checkpoint's radius at the same
severity. SR instead compares that raw radius with clean different-state
distances from the checkpoint under evaluation. The fixed-threshold results
therefore cannot be reproduced by nesting one condition inside the other.

For each task and training run, the grid contains eight training-noise maxima
and four evaluation severities. Let $x_{ij}$ denote any diagnostic or behavioral
quantity at training level $i$ and evaluation severity $j$. We remove the two
additive main effects using
\[
\widetilde{x}_{ij}
=x_{ij}-\overline{x}_{i\cdot}-\overline{x}_{\cdot j}
+\overline{x}_{\cdot\cdot}.
\]
Spearman correlations are computed over the \vTwoFactorialCellCount{} centered
cells within a single task and training run. We then average the
\vTwoRunCount{} training-run correlations within each task and weight the
\vTwoTaskCount{} tasks equally. The task ranges in
\Cref{tab:multiseverity-associations} summarize heterogeneity; they are not
confidence intervals. Double-centering leaves
$(\vTwoTrainLevelCount-1)(\vTwoEvalLevelCount-1)=\vTwoInteractionDf$
interaction degrees of freedom and imposes row- and column-sum constraints, so
the \vTwoFactorialCellCount{} cells are not treated as independent observations.

\input{\vTwoAssociationTable}
\FloatBarrier

On PushT, the exposure comparator's direct task-level correlation exceeds both
diagnostics. This is why we use it as an explicit design-metadata comparator
and do not interpret the small differences among the direct aggregate values as
diagnostic superiority.

Using stress success divided by same-checkpoint clean success instead of their
percentage-point difference gives the same qualitative result. The equal-task
mean direct correlations are $\vTwoIRRatioDirect$ for lower relative IR and
$\vTwoSRRatioDirect$ for higher SR; after accounting for the exposure ratio
they are $\vTwoIRRatioAdjusted$ and $\vTwoSRRatioAdjusted$, respectively.
Every clean-success estimate is nonzero (minimum
$\vTwoMinimumCleanSuccess\%$), but ratios can still amplify changes for weak
checkpoints; we therefore keep the percentage-point difference as the primary
outcome and treat the ratio only as a sensitivity check. These are descriptive
factorial associations. The \vTwoRowCount{} rows are not independent
replicates, and the four evaluation-severity rows for a checkpoint share the
same clean-performance estimate; therefore we do not attach row-level
$p$-values or confidence intervals.
